\documentclass[10pt,a4paper]{amsart}
\usepackage[margin=19mm]{geometry}
\usepackage{amsmath,amssymb,mathtools}
\usepackage[scr=rsfs,cal=euler]{mathalfa}
\usepackage{xfrac}
\usepackage[unicode,hidelinks,bookmarks=false]{hyperref}
\newcommand{\ie}{i.e.\ }
\newcommand{\cf}{cf.\ }
\newcommand{\resp}{respectively\ }
\newcommand{\Subsection}{\subsection}
\providecommand{\nobreakdash}{\nobreak\hbox{-}}
\setlength{\emergencystretch}{2em}
\multlinegap0pt
\usepackage[shortlabels]{enumitem}
\usepackage{amssymb}
\usepackage{mathtools}
\usepackage{xcolor}
\usepackage{tikz}
\newtheorem{lemma}{Lemma}
\title{On $k$-connected vertex-pancyclic graphs without pancyclic edges}
\author{Leyou Xu, Bo Zhou}
\date{Source: arXiv:2605.21165v1 (CC BY 4.0)}
\begin{document}
\maketitle

\noindent\emph{Source and licence.} On $k$-connected vertex-pancyclic graphs without pancyclic edges. Version arXiv:2605.21165v1 (CC BY 4.0), distributed by arXiv under the Creative Commons Attribution 4.0 International licence, http://creativecommons.org/licenses/by/4.0/, as recorded in the arXiv licence metadata for this exact version and shown on the abstract page https://arxiv.org/abs/2605.21165v1. Full licence notice: \texttt{licenses/arxiv-2605.21165-CC-BY-4.0.txt}.\par\smallskip

\noindent\textit{Editorial scope.} Counted unit: the article's lemma lem:proper-cycles (source0.tex lines 675--677). The article's complete original proof of it is delivered with it (source0.tex lines 679--727, one \texttt{proof} environment of 49 lines). No mathematics is written, repaired or re-derived: the statement and the proof are the article's own lines, in the article's own notation, and no retained line is substituted or re-flowed.  No further source-local context is needed: every cross-reference of every retained line resolves inside the two delivered spans.  Nothing was omitted from any retained span, so the package carries no omission record.  No retained line contains a citation, so the wrapper carries no bibliography and no bibliography entry.  No retained line contains a figure environment, an image-inclusion command or drawing code, so no image is delivered and the package declares no asset dependency.  Editorial formatting only, in addition: an AMS \texttt{amsart} preamble that restates the article's own declarations and macros used by the retained lines, namely the environments \texttt{lemma} on separate counters, together with the standard packages those lines need.  The retained environments print in this document as Lemma 1 in order of appearance, whereas the article prints the same items with the numbering of its own full document.  The complete native source of the article as distributed in the arXiv e-print for this exact version is bundled unchanged as \texttt{sources/201044/source0.tex}.
\begin{lemma}\label{lem:proper-cycles}
For every $q$ with $3\le q\le 8k$ and every vertex $x^{(r)}$ of $R_k$, there exists a properly colored $q$-cycle in $R_k$ containing $x^{(r)}$.
\end{lemma}

\begin{proof}
Let
\[
a=100,
\qquad
b=010,
\qquad
c=001.
\]
We first list six paths in any layer (we omit the layer number here), all starting at $z=000$ in the following table,  the color of an edge between two labels $x,y\in Q$ is the unique index $i$ such that $x+y\in P_i$. 
%The displayed table is checked directly from the definitions of $P_0,P_1,P_2$.
\[
\begin{array}{c|l|l}
\ell & \text{path }L_\ell & \text{color sequence along }L_\ell\\
\hline
3 & zc(b+c) & 1,0 \\
4 & zcb(b+c) & 1,2,1 \\
5 & zcb(b+c)(a+b+c) & 1,2,1,0\\
6 & zcba(b+c)(a+b+c) & 1,2,1,2,0 \\
7 & zcb(b+c)a(a+c)(a+b+c) & 1,2,1,2,1,0 \\
8 & zcb(b+c)a(a+c)(a+b)(a+b+c) & 1,2,1,2,1,2,1
\end{array}
\]
 In particular, for every $\ell\in\{3,4,5,6,7,8\}$:
\begin{enumerate}[label=(\alph*)]
\item $L_\ell$ is a properly colored path on $\ell$ vertices;
\item the first edge of $L_\ell$ has color $1$;
\item the last vertex of $L_\ell$ belongs to $P_2$;
\item the color of the last edge of $L_\ell$ is different from $2$.
\end{enumerate}

Now fix $q$ with $3\le q\le 8k$ and fix a vertex $x^{(r)}$ of $R_k$. Let $s=\left\lceil \frac{q}{8}\right\rceil$.
Then $s\le k$ and
\[
3s\le q\le 8s.
\]
Therefore we may write
$q=q_1+q_2+\cdots+q_s$
with each $q_j\in\{3,\dots,8\}$. %: start with $q_j=3$ for all $j$ and distribute the remaining $q-3s\le 5s$ units among the $s$ terms.

Choose distinct layers $r_1,\ldots,r_s$ with $r_1=r$. In layer $r_j$, place a `translated' copy of the path $L_{q_j}$, obtained by adding $x$ to every label. We denote the resulting path by $L_j^*$. Thus $L_j^*$ starts at $x^{(r_j)}$ and ends at $(x+h_j)^{(r_j)}$ for some $h_j\in P_2$.

For $j=1,\ldots,s$, note that $(x+h_j)^{(r_j)}x^{(r_{j+1})}\in E(R_k)$, where the index $j$ takes modulo $s$, and as $h_j\in P_2$, the edge has color $2$.
Let \[
C=x^{r}L_1^*(x+h_1)^{(r_1)}x^{(r_{2})}L_2^*\cdots L_s^*(x+h_s)^{(r_s)}x^{(r)}.
\]
Then $C$ is a cycle of length $q_1+q_2+\cdots+q_s=q$. 
For $j=1,\dots,s$, each $L_j^*$ is properly colored, and neither the first edge nor the last edge has color $2$. Hence $C$ is properly colored, as desired.
\end{proof}

\end{document}
